% !TeX root = ../colloquium-2.tex

% Формулировки критериев выпуклости в терминах первой и второй производных.
% Формулировка теоремы о взаимном расположении выпуклой функции и касательной. 
% Неравенство Йенсена.

\phantomsection\label{sec_Definition19}
\begin{definition}[
    Формулировки критериев выпуклости в терминах первой и второй производных.
    Формулировка теоремы о взаимном расположении выпуклой функции и касательной. 
    Неравенство Йенсена.
]
    \subsection*{Формулировки критериев выпуклости в терминах первой и второй производных.}
    \begin{note}[purple]
        Пусть функция $f$ дифференцируема на интервале $(a, \; b)$. Тогда выпуклость $f$ 
        на этом интервале равносильна тому, что функция $f'$ не убывает на $(a, \; b)$.
    \end{note}

    \subsection*{Формулировка теоремы о взаимном расположении выпуклой функции и касательной.}
    \begin{note}[purple]
        Пусть функция $f$ дважды дифференцируема на интервале $(a, \; b)$. Тогда выпуклость $f$
        на этом интервале равносильна тому, что функция $f''$  неотрицательна на $(a, \; b)$. 
    \end{note}

    \subsection*{Неравенство Йенсена.}
    \begin{note}[purple]
        Пусть функция $f$ выпукла на интервале $(a, \; b)$. Тогда при любых 
        $x_1, \; x_2, \; ..., \; x_n \in (a, \; b)$ и всех таких 
        $\lambda_1 \geq 0, \; \lambda_2 \geq 0, \; ..., \; \lambda_n \geq 0$, что
        $\lambda_1 + ... + \lambda_n = 1$, выполняется неравенство
        $f(\lambda_1 x_1 + ... + \lambda_n x_n) \leq \lambda_1 f(x_1) + ... + \lambda_n f(x_n)$.
    \end{note}
\end{definition}